\medskip
\textbf{Partie A}

\medskip
\begin{enumerate}
	\item 
	\begin{enumerate}
		\item L'arbre complété avec les valeurs disponibles :
		\begin{center}
				\pstree[treemode=R,nodesepA=0pt,nodesepB=2.5pt,treesep=1cm,levelsep=2.5cm]{\TR{}}
{\pstree{\TR{$V$~}\taput{$\frac23$}}
{\TR{$R$}\taput{$\frac{1}{50}$}
\TR{$\overline{R}$}\tbput{$\frac{49}{50}$}
}
\pstree{\TR{$\overline{V}$~}\tbput{$\frac13$}}
{\TR{$R$}\taput{$\frac{1}{10}$}
	\TR{$\overline{R}$}\tbput{$\frac{9}{10}$}
}
}
\end{center}
		\item D'après la formule des probabilités totales :
		
		$p(R)=p(R\cap V)+ p\left(R \cap \overline{V}\right) = p_V(R)\times p(V) + p_{\overline{V}}(R)\times p\left(\overline{V}\right) = \dfrac{1}{50}\times \dfrac23+\dfrac{1}{10}\times \dfrac13 = \dfrac{7}{150}$
		\item On cherche la probabilité : $p_R(V)$.
		
D'après la formule de Bayes, $p_R(V)=\dfrac{p(R\cap V)}{p(R)}=\dfrac{\dfrac{1}{50}\times \dfrac23}{\dfrac{7}{150}}=\dfrac{2}{7}$.
	\end{enumerate}
	\item 
	\begin{enumerate}
		\item Il s'agit là d'un schéma de Bernoulli : la répétition de 20 expériences aléatoires n'ayant que deux issues, identiques et indépendantes entre elles. $X$ suit donc une loi binomiale de paramètres $n = 20$ et $p = \dfrac23$ : $X \sim \mathcal{B}\left(20~;~\dfrac23\right)$
		\item $p(X=10)=\dbinom{20}{10}\times \left(\dfrac23\right)^5 \times \left(1-\dfrac23\right)^{20-10}\approx 0,054$.
		\item $p(X\geqslant 10)=1-p(X<10)=1-p(X\leqslant 9)\approx 0,962$
		\item Calculons l'espérance $E(X)$ de cette loi binomiale: $E(X)=n\times p=20 \times \dfrac23=\dfrac{40}{3}\approx 13,333$

En moyenne, il se rendra à la gare en vélo 13,33 jours par mois soit 13 jours à l'unité près.
	\end{enumerate}
	\item Calculons l'espérance de la loi de probabilité de $T$ :
	
	$\begin{gathered}E(t)=0,14\times 10+0,13 \times 11+0,13\times12+0,12\times13+0,12\times14+0,11\times15+0,10\times16+0,08\times17+0,07\times18\\
	=13,5\end{gathered}$
	
	Il mettra en moyenne 13 minutes et demie  pour se rendre à la gare en voiture.
\end{enumerate}

\pagebreak
